\chapter{Wnioski}
\label{cha:wnioski}

\section{Podsumowanie realizacji}
\label{sec:podsumowanie}

W ramach projektu zrealizowano kompletny pipeline klasyfikacji znaków dłoni:
od przygotowania danych tablicowych, przez uczenie i ewaluację modeli,
po inferencję z kamery w czasie rzeczywistym. Implementacja zachowuje
spójny format wejścia (28$\times$28, 784 cechy), co pozwala stosować
różne klasyfikatory bez zmian interfejsu danych.

\begin{enumerate}
      \item \textbf{CNN przewyższa metody klasyczne} na zadaniu klasyfikacji
            znaków ASL: osiąga 92\% dokładności wobec 82\% dla SVC i Lasu
            Losowego. Architektura konwolucyjna jest naturalnie dostosowana
            do danych obrazowych.
      \item \textbf{SVC i Las Losowy osiągają identyczną dokładność} (82\%)
            pomimo różnych strategii uczenia, co sugeruje, że obie metody
            osiągnęły podobną granicę dla tej reprezentacji danych.
      \item \textbf{Jakość inferencji kamerowej zależy od warunków zewnętrznych}
            --- oświetlenia, tła i pozycji dłoni --- bardziej niż od wyboru
            klasyfikatora.
      \item \textbf{MediaPipe odgrywa kluczową rolę} w przygotowaniu wejścia
            modelu; bez etapu detekcji dłoni system nie nadaje się do praktycznego
            zastosowania.
      \item \textbf{Podejście pikselowe ma istotne ograniczenia praktyczne}.
            Reprezentacja oparta na punktach kluczowych mogłaby znacząco
            poprawić odporność systemu na zmienne warunki otoczenia.
      \item \textbf{Modularna architektura projektu} (rozdzielenie warstw danych,
            modeli i inferencji) ułatwia wymianę komponentów i jest dobrą
            praktyką inżynierską w projektach uczenia maszynowego.
  \end{enumerate}

\section{Kierunki dalszego rozwoju}
\label{sec:dalszy_rozwoj}

Rozszerzenia, które mogą poprawić jakość i użyteczność systemu:
\begin{itemize}
    \item dodanie modelu sekwencyjnego (np. LSTM/Transformer) dla rozpoznawania
          gestów dynamicznych,
    \item wygładzanie predykcji w oknie czasowym,
    \item augmentacja danych i rozszerzenie zbioru treningowego o rzeczywiste
          próbki z kamery,
    \item budowa modułu oceny wydajności czasu inferencji na różnych urządzeniach.
\end{itemize}
